\section*{Глава 2. Проектирование}
\addcontentsline{toc}{section}{Глава 2. Проектирование}

\setcounter{section}{2}  
\setcounter{subsection}{1} 
\setcounter{subsubsection}{0} 

\subsubsection{Теоретический анализ целевой аудитории и разработка User Stories}

Проектируемое веб-приложение ориентировано на специфическую экосистему высшего учебного заведения. 
Целевую аудиторию системы можно сегментировать на три основные группы, каждая из которых обладает
 уникальным набором мотиваций и <<болевых точек>>.

\textbf{1. Студенты (первичная аудитория):}
\begin{itemize}
    \item \textbf{Характеристика:} Молодые люди (18–23 года), активно использующие мобильные устройства.
    \item \textbf{Проблема:} Неопределенность времени ожидания, риск <<потери>> места в бумажных списках, 
 конфликты из-за несправедливой очереди.
    \item \textbf{Потребность:} Прозрачность процесса и возможность удаленного мониторинга очереди.
\end{itemize}

\textbf{2. Преподаватели (вторичная аудитория):}
\begin{itemize}
    \item \textbf{Характеристика:} Сотрудники кафедр, заинтересованные в соблюдении тайминга занятия.
    \item \textbf{Проблема:} Сложность визуального контроля большого количества студентов, необходимость 
 вручную вести журналы успеваемости и посещаемости.
    \item \textbf{Потребность:} Автоматизация рутинных операций (вызов следующего, выставление оценок) и 
 статистика по группе.
\end{itemize}

\textbf{3. Администраторы/Старосты (связующее звено):}
\begin{itemize}
    \item \textbf{Характеристика:} Студенты с дополнительными полномочиями по управлению группой.
    \item \textbf{Проблема:} Модерация списков в Google-таблицах, борьба с <<мертвыми душами>> 
 (записавшимися, но не пришедшими).
    \item \textbf{Потребность:} Инструменты оперативного редактирования очереди и верификации состава группы.
\end{itemize}

\textbf{Разработка User Stories}
Для формализации требований к функционалу системы разработаны пользовательские истории (User Stories):

\begin{itemize}
    \item \textbf{Как Студент}, я хочу записываться в очередь через смартфон, чтобы не присутствовать физически
  у кабинета всё время ожидания.
    \item \textbf{Как Студент}, я хочу получать уведомление о необходимости подтверждения присутствия, 
 чтобы система автоматически удаляла неактивных участников и моя позиция становилась актуальной.
    \item \textbf{Как Преподаватель}, я хочу видеть список сданных работ каждого студента в очереди,
  чтобы заранее понимать объем и сложность предстоящей проверки.
    \item \textbf{Как Администратор}, я хочу иметь возможность удалять или перемещать записи, 
 чтобы оперативно разрешать спорные ситуации в аудитории.
\end{itemize}

\subsubsection{Описание ролей и сценариев взаимодействия Студент--Преподаватель--Администратор}

На основе архитектурной схемы приложения выделяются три ключевые роли, взаимодействие которых 
обеспечивает замкнутый цикл управления учебным процессом.

\textbf{1. Роль <<Студент>>}

Основной сценарий взаимодействия студента с системой начинается с \textbf{авторизации} и 
выбора конкретной дисциплины.

\begin{itemize}
    \item \textbf{Запись:} Студент отправляет запрос на сервер (<<Запись в конец очереди>>), 
	после чего его данные попадают в таблицу текущей очереди.
    \item \textbf{Подтверждение (Check-in):} За 24 часа до начала занятия срабатывает логика 
	<<Таблицы готовности>>. Студент обязан нажать кнопку подтверждения. В случае отсутствия 
	действия сервер автоматически удаляет запись, реализуя механизм самоочистки очереди.
    \item \textbf{Ожидание:} Студент получает актуальное состояние очереди по протоколу WebSocket, 
	что исключает необходимость ручного обновления страницы.
\end{itemize}

\textbf{2. Роль <<Преподаватель>>}

Преподаватель взаимодействует с наиболее сложным интерфейсом, интегрированным с блоками 
\textbf{<<Таблица успеваемости>>} и \textbf{<<Таблица посещаемости>>}.

\begin{itemize}
    \item \textbf{Управление потоком:} Преподаватель инициирует вызов студента. Сервер 
	обновляет статус записи, и через WebSocket уведомление рассылается всем участникам.
    \item \textbf{Фиксация результата:} После защиты лабораторной работы преподаватель через форму 
	редактирования отправляет данные на сервер. Система одновременно обновляет <<Журнал>> (выставляя оценку) 
	и удаляет студента из текущей активной очереди, освобождая место для следующего.
\end{itemize}

\textbf{3. Роль <<Администратор (Староста)>>}

Данная роль является расширением роли студента и обладает административными правами в рамках своей учебной группы.

\begin{itemize}
    \item \textbf{Конфигурация:} Староста отвечает за создание структуры группы (добавление студентов в 
	<<Таблицу студентов>>) и привязку предметов.
    \item \textbf{Модерация:} В сценарии <<Редактирование очереди>> администратор может вручную изменять 
	порядок записей (например, если студенты договорились о взаимозамене), что сразу отражается в базе 
	данных и на клиентах других пользователей.
    \item \textbf{Контроль посещаемости:} Администратор имеет право выставлять первичные отметки о присутствии, 
	которые затем верифицируются преподавателем.
\end{itemize}

\textbf{Логика взаимодействия (Summary):}
Центральным узлом системы выступает \textbf{Сервер}, который координирует запросы от всех типов клиентов.
 Благодаря разделению на таблицы (предметы, студенты, группы, успеваемость) обеспечивается целостность данных: 
 оценка не может быть выставлена студенту, которого нет в базе группы, а запись в очередь невозможна без выбора 
 существующего предмета.


 \subsection{Разработка UML Use Case}

Для формализации функциональных требований к системе и визуализации границ проектируемого веб-приложения используются диаграммы вариантов использования (Use Case Diagrams). Они позволяют определить взаимодействие внешних актеров (ролей) с системой.

\subsubsection{Для роли «Студент»}

Студент является массовым пользователем системы. Его взаимодействие сфокусировано на управлении собственной позицией в очереди и мониторинге учебного процесса. Основные варианты использования включают регистрацию, выбор предмета, управление записью и подтверждение присутствия.

\begin{figure}[ht!]
    \centering
     \includegraphics[width=0.8\textwidth]{diagram_students}
    \caption{UML Use Case диаграмма для роли «Студент»}
    \label{fig:uc_student}
\end{figure}

\subsubsection{Для роли «Администратор» и «Преподаватель»}

Данные роли обладают расширенными правами доступа. Согласно архитектуре приложения, Администратор (Староста) 
наследует все возможности студента, добавляя функции управления составом группы. Преподаватель же является 
«контроллером» очереди, определяющим динамику занятия.

\begin{figure}[ht!]
    \centering
     \includegraphics[width=0.8\textwidth]{UML_Use_Case_prepod_admin}
    \caption{UML Use Case диаграмма для ролей «Администратор» и «Преподаватель»}
    \label{fig:uc_admin_teacher}
\end{figure}

На диаграмме студента связи \texttt{<<include>>} указывают, что запись в очередь и подтверждение
присутствия обязательно включают в себя авторизацию, что подчёркивает безопасность системы.
Функция «Управление очередью» (UC7) доступна обоим актёрам: это соответствует архитектуре, где
администратор-староста может помочь преподавателю модерировать список в аудитории.
WebSocket-мониторинг выделен как отдельный вариант использования (UC5) для роли студента,
поскольку обновление очереди в реальном времени является ключевым нефункциональным требованием системы.

\subsection{Проектирование модели данных}

Проектирование базы данных выполняется в несколько этапов: от определения высокоуровневых сущностей 
и их взаимосвязей до формирования физической структуры таблиц, оптимизированной под выбранную СУБД PostgreSQL.

\subsubsection{Проектирование ER-диаграммы для сущностей «Администратор и Преподаватель» и логической схемы БД}

На логическом уровне определяются ключевые бизнес-сущности и правила их взаимодействия. Для ролей Преподавателя
и Администратора критически важными являются связи с учебными группами, предметами и журналами успеваемости.

Логическая схема строится на принципе минимизации избыточности. Основными узлами здесь выступают сущности «Группа»,
 которая объединяет студентов, и «Предмет», к которому привязываются очереди.

База данных приложения содержит \textbf{восемь таблиц}, каждая из которых отвечает за отдельную предметную область:

\begin{enumerate}
    \item \textbf{study\_groups} — учебные группы: идентификатор, название, дата создания.
    \item \textbf{users} — пользователи всех ролей: ФИО, email, хеш пароля, роль
    (\texttt{student}/\texttt{teacher}/\texttt{admin}), внешний ключ на группу.
    \item \textbf{subjects} — дисциплины: название, внешние ключи на преподавателя и группу,
    а также два поля, определяющих нормативы сдачи:
    \begin{itemize}
        \item \texttt{lab\_count} — количество лабораторных работ по дисциплине (1–30);
        \item \texttt{minutes\_per\_lab} — среднее время сдачи одной работы в минутах (1–120).
        Значение \texttt{minutes\_per\_lab} подставляется по умолчанию в форму записи студента
        и используется преподавателем для оценки загруженности очереди.
    \end{itemize}
    \item \textbf{queue\_sessions} — конкретные очереди на занятие: привязка к дисциплине и группе,
    дата и время начала, флаг активности.
    \item \textbf{queue\_entries} — записи студентов в очередь (подробнее в §2.3.2).
    \item \textbf{grade\_records} — журнал успеваемости: связь студент–дисциплина–номер ЛР, оценка,
    дата сдачи и ссылка на проверяющего преподавателя.
    \item \textbf{notifications} — уведомления: заголовок, текст сообщения, флаг прочтения
    (\texttt{is\_read}) и внешний ключ на пользователя-получателя.
    \item \textbf{work\_uploads} — загруженные студентами файлы работ: безопасное имя файла на диске,
    оригинальное имя, дисциплина, номер ЛР.
\end{enumerate}

\begin{figure}[ht!]
    \centering
    \includegraphics[width=0.8\textwidth]{Logical_ER_diagram_of_relationships_for_administrative_roles}
    \caption{Логическая ER-диаграмма связей для административных ролей}
    \label{fig:er_admin_teacher}
\end{figure}

\subsubsection{Проектирование ER-диаграммы для сущности «Студент» и физической схемы БД}

Физическая схема БД описывает конкретную реализацию в PostgreSQL, включая типы данных, первичные и внешние
ключи (PK/FK), а также ограничения (Constraints). Схема охватывает все восемь таблиц приложения,
сформированных через декларативные модели SQLAlchemy.

Центральной таблицей является \texttt{users}, хранящая всех пользователей системы (студентов, преподавателей
и администраторов) с разграничением через поле \texttt{role}. Студент взаимодействует главным образом
с тремя таблицами: \texttt{subjects} (список дисциплин), \texttt{queue\_sessions} (открытые сессии) и
\texttt{queue\_entries} (его собственные записи в очередь). Записи в \texttt{queue\_entries} содержат поле
\texttt{position} и временну\'{ю} метку \texttt{joined\_at} для реализации алгоритма FIFO, а также поле
\texttt{status} со значениями \texttt{waiting}/\texttt{checking}/\texttt{passed}/\texttt{cancelled}.

\begin{figure}[ht!]
    \centering
   \includegraphics[width=0.8\textwidth]{Physical_data_schema_Student_entity}
    \caption{Физическая схема базы данных приложения (все 8 таблиц)}
    \label{fig:student_physical}
\end{figure}

Полный перечень связей между таблицами с указанием кардинальности и поведения при удалении
приведён в таблице~\ref{tab:db_relations}.

\begin{table}[ht!]
\centering
\caption{Связи между таблицами базы данных}
\label{tab:db_relations}
\small
\begin{tabularx}{\textwidth}{p{3.2cm}|p{3.2cm}|p{2.5cm}|p{1.8cm}|p{3cm}}
\hline
\textbf{Таблица (FK)} & \textbf{Ссылается на} & \textbf{Поле FK} & \textbf{Тип} & \textbf{ON DELETE} \\
\hline
users & study\_groups & group\_id & M:1 (opt.) & SET NULL \\
\hline
subjects & users (teacher) & teacher\_id & M:1 (opt.) & SET NULL \\
\hline
queue\_sessions & subjects & subject\_id & M:1 & CASCADE \\
\hline
queue\_sessions & users (teacher) & teacher\_id & M:1 & CASCADE \\
\hline
queue\_entries & queue\_sessions & session\_id & M:1 & CASCADE \\
\hline
queue\_entries & users (student) & student\_id & M:1 & CASCADE \\
\hline
grade\_records & users (student) & student\_id & M:1 & CASCADE \\
\hline
grade\_records & subjects & subject\_id & M:1 & CASCADE \\
\hline
grade\_records & queue\_sessions & session\_id & M:1 (opt.) & SET NULL \\
\hline
grade\_records & users (teacher) & teacher\_id & M:1 (opt.) & SET NULL \\
\hline
work\_uploads & users (student) & student\_id & M:1 & CASCADE \\
\hline
work\_uploads & subjects & subject\_id & M:1 & CASCADE \\
\hline
work\_uploads & queue\_sessions & session\_id & M:1 (opt.) & SET NULL \\
\hline
notifications & users & user\_id & M:1 & CASCADE \\
\hline
\end{tabularx}
\end{table}

Обозначение «opt.» (optional) означает, что внешний ключ допускает значение \texttt{NULL} —
например, дисциплина может существовать без назначенного преподавателя (\texttt{teacher\_id}).
Связи типа \texttt{ON DELETE CASCADE} обеспечивают автоматическое удаление зависимых записей;
\texttt{ON DELETE SET NULL} — обнуление ссылки без удаления строки.

\textbf{Обоснование структуры физической модели:}

\begin{enumerate}
    \item \textbf{Тип данных SERIAL:} Используется для автоматической инкрементации идентификаторов,
 что обеспечивает высокую скорость вставки новых записей в очередь.
    \item \textbf{Поле created\_at:} Является критическим для алгоритма FIFO. Сортировка очереди на сервере
 Flask выполняется по полям \texttt{position} и \texttt{created\_at}.
    \item \textbf{Поля \texttt{lab\_number} и \texttt{work\_count}:} Позволяют студенту указать,
 какую именно лабораторную работу он сдаёт и сколько работ планирует сдать за одно посещение.
    \item \textbf{Поле \texttt{estimated\_minutes}:} Студент указывает предполагаемое время сдачи;
 по умолчанию подставляется значение \texttt{minutes\_per\_lab} из таблицы \texttt{subjects}.
 Это позволяет преподавателю оценить общую продолжительность очереди.
    \item \textbf{Поле \texttt{status}:} Хранит состояние записи (\texttt{waiting}, \texttt{checking},
 \texttt{passed}, \texttt{cancelled}). Только записи со статусом \texttt{waiting} и \texttt{checking}
 отображаются в активной очереди.
    \item \textbf{Таблица notifications:} Вынесена отдельно для хранения уведомлений независимо
 от текущего состояния очереди; флаг \texttt{is\_read} позволяет реализовать счётчик непрочитанных.
    \item \textbf{Связи ON DELETE CASCADE:} Использование внешних ключей с каскадным удалением обеспечивает
 целостность данных — при удалении студента из системы все его записи в очередях и оценки удаляются автоматически.
\end{enumerate}

\paragraph{Пояснение по логике подтверждения:}
В таблице \texttt{queue\_entries} поле \texttt{is\_confirmed} по умолчанию имеет значение \texttt{false}. 
Фоновый процесс (task) на сервере или логика контроллера проверяет это поле за определенное время до начала
 занятия. Если к моменту «отсечки» значение остается \texttt{false}, запись переходит в статус 
 \texttt{cancelled} или удаляется физически, что мгновенно отражается у всех пользователей через WebSockets.


\subsection{Разработка макетов Wireframes}

Проектирование интерфейса (UI) системы «Управление очередями» велось с учетом принципа «Mobile First», 
так как основная часть пользователей (студенты) взаимодействует с системой через смартфоны в условиях учебных 
аудиторий.

\subsubsection{Разработка макета Wireframe для очереди и Студента}

Интерфейс студента спроектирован таким образом, чтобы минимизировать количество действий для совершения ключевых 
операций: записи в очередь и подтверждения присутствия.

\textbf{Основные элементы макета:}

\begin{enumerate}
    \item \textbf{Информационная панель:} Отображение текущей дисциплины и статуса сессии.
    \item \textbf{Виджет подтверждения:} Яркая кнопка «Подтвердить присутствие», которая появляется только 
 в активное окно (за 24 часа).
    \item \textbf{Динамический список:} Таблица со столбцами «№», «ФИО», «Статус». Собственная позиция студента
  подсвечивается контрастным цветом.
    \item \textbf{Кнопка действия:} Крупная кнопка «Встать в очередь» (или «Выйти из очереди»), расположенная в 
 нижней части экрана для удобства нажатия одной рукой.
\end{enumerate}

\textbf{Mакет интерфейса студента:}

\begin{figure}[ht!]
    \centering
   \includegraphics[width=0.8\textwidth]{example_pril_ochered}
    \caption{Пример макета интерфеса для студента}
    \label{fig:ochered_example}
\end{figure}
\subsubsection{Разработка макета Wireframe для Преподавателя и администратора}

Интерфейс для ролей с расширенными правами ориентирован на управление и быстрый ввод данных. 
В отличие от студенческого макета, здесь добавлены элементы управления каждой строкой очереди.

\textbf{Основные элементы макета:}

\begin{enumerate}
    \item \textbf{Панель управления сессией:} Кнопки «Начать прием», «Пауза», «Завершить».
    \item \textbf{Интерактивная очередь:} Каждая строка содержит кнопки «Вызвать», «Удалить», «Передвинуть».
    \item \textbf{Форма быстрого ввода (Modal):} При нажатии на студента открывается окно для 
	выставления оценки за конкретную лабораторную работу.
    \item \textbf{Вкладка «Журнал»:} Переключатель на просмотр общей ведомости группы 241-326.
\end{enumerate}

\textbf{Mакет интерфейса Преподавателя:}

\begin{figure}[ht!]
    \centering
   \includegraphics[width=0.8\textwidth]{examle_pril_admin_prepod}
    \caption{Пример макета интерфеса для преподавателя и администратора}
    \label{fig:ochered_admin_prepod}
\end{figure}


\subsubsection{Разработка макета Wireframe для Администратора (Старосты)}

Роль Администратора (Старосты) является расширением роли студента и предоставляет доступ к
инструментам управления составом группы. Интерфейс ориентирован на административные операции:
добавление пользователей, создание учебных групп и назначение дисциплин.

\textbf{Основные элементы макета:}

\begin{enumerate}
    \item \textbf{Панель управления группами:} Список существующих групп с кнопкой создания новой,
    отображение количества студентов в каждой группе.
    \item \textbf{Управление пользователями:} Таблица всех пользователей системы с фильтрацией по роли
    и группе; возможность смены роли и привязки к группе без перерегистрации.
    \item \textbf{Форма быстрого создания:} Компактная форма добавления пользователя прямо
    на странице без перехода на отдельный маршрут.
    \item \textbf{Доступ к очередям:} Наряду с административными инструментами, Администратор
    видит те же очереди, что и Преподаватель, и может модерировать порядок записей.
\end{enumerate}

\textbf{Макет интерфейса Администратора (Старосты):}

\begin{figure}[ht!]
    \centering
    \includegraphics[width=0.85\textwidth]{wireframe_admin}
    \caption{Пример интерфейса для роли «Администратор (Староста)»}
    \label{fig:wireframe_admin}
\end{figure}

\textbf{Технологическое обоснование дизайна:}

\begin{itemize}
    \item \textbf{Сетка (Grid):} Для верстки будет использоваться CSS Flexbox, что позволит макетам 
	одинаково корректно отображаться как на десктопе (у преподавателя), так и на узких экранах смартфонов.
    \item \textbf{Цветовая индикация:} Статусы (Сдает, Ожидает, Подтвержден) будут иметь цветовую кодировку 
	(зеленый, желтый, синий), что позволит пользователю мгновенно считывать ситуацию без чтения текста.
    \item \textbf{WebSockets:} При нажатии преподавателем кнопки «Вызвать», у студента на макете мгновенно
	 обновится статус, а позиция «ВЫ» переместится в блок «ТЕКУЩИЙ».
\end{itemize}